\section{From one-field purchases to the public comparison grid}
\label{ex-public-grill-producer}

This section isolates the quantitative steps on the selected finite-volume
route to the comparison grid. It proves the finite-frame price estimate, the
radial decay once the trajectory bounds hold, and the realification and
four-term assembly steps. The local approximation and separated-pair
receipts are substantive upstream inputs; they are recorded explicitly at
the end rather than treated as consequences of the assembly.

Write $a_j$ for the physical spacing at depth $j$, $c_j$ for the number of
columns there, and $\sigma_j$ for its running scale. For a coefficient $L$
of grade $d$, let $\|L\|_1$ be its coefficient norm in the finite frame.
The price appearing in the construction is
\[
  K_j(L)=2\exp(\ell_j)\|L\|_1,
\]
where $\ell_j$ is the accumulated column logarithm. The rooted price of a
block of side at depth $m$, compared with a fine observable at depth $j$, is
\[
  S_{m,j}=\exp\!\left(\left(\frac{a_m}{a_j}\right)^4
                        \mathcal T_j+5\right),
\]
with $\mathcal T_j$ the vacuum treasury. This is the four-dimensional
marked-copy estimate proved in Section~\ref{ex-centered-sail}.

\begin{lemma}[Finite-frame recovery from approximate coordinate atoms]
\label{lem:ex-public-frame-recovery}
Let $V$ be a finite-dimensional complex coefficient space with coordinate
norm $\|\cdot\|_1$ and basis $(e_s)_{s\in I}$. Suppose a family
$(y_s)_{s\in I}$ satisfies $\|e_s-y_s\|_1\leq1/4$ for every $s$.
Define $A v=\sum_{s\in I}v_s y_s$. Then $A$ is invertible and
$\|A^{-1}z\|_1\leq2\|z\|_1$ for all $z\in V$.
\end{lemma}

\begin{proof}
For $v=\sum_s v_s e_s$, the triangle inequality gives
\[
 \|v-Av\|_1
   =\left\|\sum_s v_s(e_s-y_s)\right\|_1
   \leq\sum_s|v_s|\|e_s-y_s\|_1
   \leq\tfrac14\|v\|_1.
\]
Consequently $\|Av\|_1\geq3\|v\|_1/4$. Thus $A$ is injective and,
by finite dimensionality, surjective. The displayed inequality gives
$\|A^{-1}z\|_1\leq4\|z\|_1/3\leq2\|z\|_1$.
\end{proof}

\begin{lemma}[Linear assembly preserves the local price]
\label{lem:ex-public-origin-sum}
Fix a finite coefficient frame and a compact local region $W$. Suppose the
complex basis approximants $Z_s$ satisfy
\[
 \sup_x|Z_s(x)|\leq A_0,\qquad Z_s\text{ is supported in }W,
 \qquad \|Z_s-f_s\|_{\mathrm{root},m}\leq\epsilon_0
 \quad(s\in I),
\]
where $f_s$ are the finite-volume insertions corresponding to the frame
columns. Let $\mathcal C$ denote the centered chamber-interior test class;
assume $Z_s,f_s,Z_s-f_s\in\mathcal C$ and, for finite combinations,
\[
 \rho\!\left(\sum_s t_s H_s\right)
   \leq\sum_s|t_s|\rho(H_s)\qquad(H_s\in\mathcal C),
\]
where $\rho$ is the rooted test. If $L=\sum_s t_s f_s$ with
$\sum_s|t_s|\leq K$, then
\[
 Z=\sum_s t_s Z_s
\]
is supported in $W$, belongs to $\mathcal C$, satisfies
$\sup_x|Z(x)|\leq A_0K$, and
\[
                 \|Z-L\|_{\mathrm{root},m}\leq K\epsilon_0.
\]
\end{lemma}

\begin{proof}
Support is preserved under finite linear combinations. For every $x$,
$|Z(x)|\leq\sum_s|t_s||Z_s(x)|\leq A_0K$. Also
$Z-L=\sum_s t_s(Z_s-f_s)$, so subadditivity and homogeneity give the root
bound. The same closure assumptions give $Z-L\in\mathcal C$. No independence
assumption is used: all rooted tests are evaluated in the same finite-volume
Wilson law.
\end{proof}

\begin{proposition}[A finite-depth real approximation receipt]
\label{prop:ex-public-origin}
Fix grade $d$, a star-real coefficient $L$, a block depth $m$, a station
$J\leq N$, and a fine depth $k<N$ with $C.ref(d)\leq m\leq k$. Assume the
following selected-history data at these depths. The rooted quantity below
is the source root test for the same finite-volume law and centered chamber;
we assume its triangle inequality, homogeneity, and invariance under complex
conjugation, as in the selected root-test receipts.
\begin{enumerate}
\item each coordinate coefficient has a complex local approximation in the
      window of radius $R_d+4$, with supremum bound $A_0$ and rooted error at
      most $S_{J,m}/w$;
\item the finite-volume insertion columns span the requested insertion $L$,
      with coefficient sum at most $K_k(L)$, and the resulting complex sum
      $Z$ lies in the centered chamber interior; that interior is preserved
      by the real-part projection used below;
\item moving each local approximation from the fine centered box back to the
      block of depth $m$ preserves its support inside the block and its rooted
      error bound;
\item the scales fit: the local window embeds in the block, the station box
      is open at depth $m$, and
\[
 w\,\|\operatorname{alto}_{m,d}\|
       \frac{4(1+c_k)}{\operatorname{gain}(k,m)}
       \leq \frac{\zeta_{d,m}}2,
 \qquad w>0,\quad \zeta_{d,m}>0.
\]
\end{enumerate}
Then there is a real field $B_N$ supported in the centered block, uniformly
bounded by $A_0K_k(L)$, whose rooted distance from the insertion $L$ is at
most
\[
                       K_k(L)\,\frac{S_{J,m}}w.
\]
The same bound holds at every period once the station and scale inequalities
hold there.
\end{proposition}

\begin{proof}
Use the representation and coefficient bound in item~2. Transport the
complex coordinate approximants to the
common block and form their complex linear combination $Z$. The finite-sum
lemma gives $|Z|\leq A_0K_k(L)$ and rooted error at most
$K_k(L)S_{J,m}/w$. The exact selected-family reality identity says that the
target insertion $U$ is real when $L$ is star-real. Put $B=\operatorname{Re}Z$.
Then $|B|\leq|Z|$, and $\operatorname{paint}(B)$ has the same local support.
The selected realification lemma, namely chamber-interior compatibility for
averaging a field with its conjugate, gives its chamber interior. Since
$U=\overline U$,
\[
 \operatorname{paint}(B)-U
   =\tfrac12\bigl((Z-U)+\overline{(Z-U)}\bigr).
\]
Conjugation preserves the rooted test, and its triangle inequality and
homogeneity therefore give the same rooted error bound for $B$. This is the
realification argument of the selected route, not an assumption that the
complex coordinate approximants are themselves real. The support condition
and scale inequality ensure that all summands are legal tests in one block.
Uniformity in period follows from the common station and aligned physical
block in the hypotheses.
\end{proof}

\begin{lemma}[Replacing a finite local word]
\label{lem:ex-public-replace-word}
Let $\mathcal I$ be a finite index set, $\mathcal T$ a finite tiling, and
$\tau_a$ its copy maps, indexed by the cells $a\in\operatorname{Cell}(\mathcal
T)$. Assume each copy map is additive. Let $\mu$ be a probability law and
$\rho$ the same nonnegative rooted test on all functions considered. Let
$a:\mathcal I\to\operatorname{Cell}(\mathcal T)$ be injective. Suppose
$X_i,Y_i$ lie in a linear local test class that is closed under cell copies,
and suppose all finite products below are integrable.
\[
 \rho(X_i)\leq M,\qquad \rho(Y_i)\leq M,\qquad
 \rho(Y_i-X_i)\leq\delta
 \quad(i\in\mathcal I),
\]
where $M,\delta\geq0$. Assume the planted-product estimate
\[
 \left|\int\prod_{i\in\mathcal I}\tau_{a(i)}(Z_i)\,d\mu\right|
       \leq\prod_{i\in\mathcal I}\rho(Z_i)
\]
for every family in this local test class. If $|\mathcal I|=n\geq1$, then
\[
 \left|\int\prod_{i\in\mathcal I}\tau_{a(i)}(Y_i)\,d\mu
       -\int\prod_{i\in\mathcal I}\tau_{a(i)}(X_i)\,d\mu\right|
       \leq n\delta M^{n-1}.
\]
For $n=0$ the two empty products agree and the difference is zero.
\end{lemma}

\begin{proof}
Order the finite index set as $i_1,\ldots,i_n$, and define $Z^{(r)}$ to use
$Y_{i_s}$ for $s\leq r$ and $X_{i_s}$ for $s>r$. Then the difference of the
two products is the sum, for $r=1,\ldots,n$, of the product with
$\tau_{a(i_r)}(Y_{i_r}-X_{i_r})$ in position $r$ and the hybrid
factors in all other positions. The copy operation is additive, including
its reflected cases, so this is an exact telescoping identity. Each hybrid
family remains in the local test class. The planted-product estimate bounds the
absolute value of the $r$th integral by $\delta M^{n-1}$: the changed factor
costs at most $\delta$ and each of the other $n-1$ factors costs at most
$M$. Summing these $n$ bounds proves the claim.
\end{proof}

\begin{remark}[The selected pair receipt]
\label{rem:ex-public-layout-receipt}
This is the quantitative replacement step used to turn centered single-word
receipts into whole-word product receipts. Its local inputs are: for each
copy and each mesh/period in the
fused late set, the exact insertion $X_i$ and real prepared insertion $Y_i$
lie in the centered chamber-interior test class, both rooted costs are at most a common
$M$, their difference costs at most $\operatorname{costHint}(k)$, and the
prepared field has a finite supremum bound. The finite-dimensional spin
frame is absorbed by summing its centered and purchase budgets. For an
injective list of slots, the measure's planted inequality is applied to each
telescoping term above. Thus the word-product discrepancy is bounded by
$n\,\operatorname{costHint}(k)M^{n-1}$, uniformly over the finite grid
positions; concatenating two finite lists gives the four pair receipts used
in the reflection-difference assembly. The selected proof still needs the
origin receipts that feed the centered single-word estimates, but the
many-factor replacement estimate itself is proved here.
\end{remark}

\begin{lemma}[The radial choice makes the coefficient price affordable]
\label{lem:ex-public-radial-choice}
Fix $d,m,J$ and $\varepsilon>0$. Suppose the scale schedule satisfies
\[
\sigma_j\longrightarrow+\infty,\qquad
\ell_j=o(\sigma_j),\qquad \log(1+c_j)=o(\sigma_j),\qquad
 \operatorname{gain}(j,m)^{-1}\leq C_m e^{-\gamma(\sigma_j-\sigma_m)}
\]
for some $\gamma>0$. Assume $\|\operatorname{alto}_{m,d}\|>0$,
$\zeta_{d,m}>0$, and $\operatorname{gain}(j,m)>0$. Let $w_j$ be the positive allowance
\[
 w_j=\frac{\zeta_{d,m}\operatorname{gain}(j,m)}
              {8\|\operatorname{alto}_{m,d}\|(1+c_j)}.
\]
If $\|L\|_1$ and $\zeta_{d,m}^{-1}$ are fixed, then for all sufficiently
large $j$,
\[
              K_j(L)\,\frac{S_{J,m}}{w_j}\leq\varepsilon.
\]
\end{lemma}

\begin{proof}
For fixed $d,m,J,L$, all factors except $\exp(\ell_j)$,
$1+c_j$, and $\operatorname{gain}(j,m)^{-1}$ are constant in $j$; the
remaining treasury ratio in $S_{J,m}$ is fixed at the chosen $m,J$. By the
first two hypotheses, for any $\eta>0$ both $\exp(\ell_j)$ and $1+c_j$ are
eventually at most $\exp(\eta\sigma_j)$. Choose $\eta<\gamma/3$. The gain
factor is bounded by a constant times $e^{-\gamma\sigma_j}$, after absorbing
$e^{\gamma\sigma_m}$ into that constant. The product is thus bounded by a
constant times $e^{-(\gamma-2\eta)\sigma_j}$, which tends to zero. This gives
the claimed threshold. The source obtains the first two subpower estimates
from the accumulated column logarithm and column-count bounds; the gain decay
is the scale-flight estimate.
\end{proof}

\begin{theorem}[Conditional assembly of the public comparison grid]
\label{thm:ex-public-grill}
Fix a finite list of star-real tokens $v_i$ satisfying the source plainness
condition, and a generic grid of pairwise separated positions $x_i$. Write
$H_i$ and $R_i$ for the exact unreflected and reflected whole-word insertions
at common volume, period, and translation; reflection acts on both labels
and positions. For each approximation index $k$, suppose there are real
measurable approximants $F_{k,N,p,t,i}$ and $G_{k,N,p,t,i}$ to the
unreflected and reflected insertions. Assume:
\begin{enumerate}
\item $A(k)\geq0$, $\delta(k)\geq0$, and $\delta(k)\to0$. For each fixed
      $k,t,i$, eventually in the fused mesh/period filter $f.\mathrm{fuse}$,
      $F_{k,N,p,t,i}$ and $G_{k,N,p,t,i}$ are measurable and bounded in
      absolute value by the common reserve $A(k)$.
\item For each $k$ and each $i<j$, there is a translation threshold such
      that for every $t$ beyond it, eventually in $f.\mathrm{late}$ in
      $(N,p)$, the four approximating products and their target products are
      integrable and their integrals differ by at most $\delta(k)$:
      \begin{align*}
       \int F_iF_j&\leftrightarrow\int\operatorname{Re}(H_iH_j),&
       \int F_iG_j&\leftrightarrow\int\operatorname{Re}(H_iR_j),\\
       \int G_iF_j&\leftrightarrow\int\operatorname{Re}(R_iH_j),&
       \int G_iG_j&\leftrightarrow\int\operatorname{Re}(R_iR_j).
      \end{align*}
      The arrows mean that the absolute difference of the two real numbers is
      at most $\delta(k)$.
\end{enumerate}
Then the real plots $P_{k,N,p,t,i}=F_{k,N,p,t,i}-G_{k,N,p,t,i}$,
with reserve $2A(k)$ and error schedule $4\delta(k)$, define a covering
satisfying the exact Works quantifiers: controls hold eventually in
$f.\mathrm{fuse}$ for each fixed $k,t,i$; and for every $k$ and $i<j$, its
pair product differs from the bill
\[
 \int\operatorname{Re}\bigl[(H_i-R_i)(H_j-R_j)\bigr]
\]
by at most $4\delta(k)$ eventually in $f.\mathrm{late}$ for every
sufficiently large translation $t$. Thus $k$ is the approximation index
whose error tends to zero; it is not the extraction index. This is a
conditional finite assembly statement, not a proof of the four local pair
receipts.
\end{theorem}

\begin{proof}
Put $P_i=F_i-G_i$. The bound $|P_i|\leq |F_i|+|G_i|\leq2A(k)$ proves the
reserve assertion. The identity
\[
 P_iP_j=F_iF_j-F_iG_j-G_iF_j+G_iG_j
\]
is an exact identity of real functions. The target bill has the same
four-term expansion after taking real parts of the exact complex products
$H_iH_j,H_iR_j,R_iH_j,R_iR_j$. By the four assumed receipts and the triangle
inequality, the difference of the two integrals is at most $4\delta(k)$.
The common late set exists because there are finitely many controls and four
receipts per pair. The translation threshold is chosen first, followed by
the mesh/period late set. This is precisely the order of quantifiers in
Works.
\end{proof}

\begin{remark}[Selected-route boundary]
The public endpoint reaches this construction through the selected run's
history records, then uses the finite-dimensional insertion recovery, a
finite real sum, a radial choice of the intermediate depth, and the
four-term reflection expansion. The arguments above prove the coefficient
recovery, finite-sum estimates, radial decay under the displayed scale
bounds, and the final four-term assembly. The selected-history receipts that
remain inputs are: the real local coordinate approximations, their rooted
one-field errors, the back-transport support and price estimate, and the
separated pair integral estimates. These are the substantive content of the
history lemmas producing the local approximants and their served pair bounds;
this section does not treat their names alone as proofs. The public theorem
also separately requires same-history full-mesh convergence for positioned
correlations. This chapter supplies the real grid estimate on its own
hypotheses and does not identify it with that convergence theorem.
\end{remark}

\section{The selected one-scale linear interpolation}
\label{ex-selected-linear-interpolation}

We next expand the finite construction that turns the one-column coupons into
one coherent coarser voice. It is essential that the coupons be combined
before the coefficient is chosen: choosing separately for each label would
not give a single linear voice, and the next scale's cone estimate would not
apply. Let $\mathcal T_m$ denote the actual fine-to-coarse density map,
including its physical block-volume factor. It is integration against the
selected transition kernel, so its linearity on the displayed integrable
columns follows from the probability-kernel identity.

Let $F$ be the finite coefficient frame, $V$ a voice at the fine scale, and
write $I$ for its invertible coefficient readout. Assume the heard identities
hold: the played coefficient has $x$-coordinate $I(v)$, its density is a
complex-linear function of $v$, the density obeys the stated motion
covariance, and its charge is bounded by $s\|\operatorname{core}I(v)\|_1$,
where $s$ is the fine column count. For the profile conclusions below assume
the coefficient law is sound under local reorientation, the forced profiles
transform covariantly, and the local evaluation is nonzero on every
configuration. Density covariance alone does not imply covariance of the
exponential entry. Let
$L$ be the block scale and $T$ the selected heading. Set
\[
 D(z)=\|\operatorname{vowel}_L(\operatorname{core}z)\|_1,
 \quad c=C+\eta s,\quad \epsilon=\eta(1+s),
 \quad R_F=\operatorname{orbitCost}(F).
\]
The frame has basis $(e_j)_{j\in J}$ and finite motion set
$\mathcal M=\mathrm{Signed}\times\{0,1\}$; the second coordinate is the
independent conjugation flag. Write $\sigma_g(c)=c$ or $\bar c$ according to
that flag, and $g z$ for the corresponding coefficient action. The exact
repainting coefficient is
\[
 \chi_{g,j}(z)=|\mathcal M|^{-1}\sigma_g((g^{-1}z)_j),
 \qquad R_F=1+\sum_{g\in\mathcal M}\sum_{j\in J}\|g(e_j)\|_1.
\]
It is complex-linear: applying $g^{-1}$ to $cz$ contributes
$\sigma_{g^{-1}}(c)$, and the outer $\sigma_g$ changes it back to $c$.
For fixed $g$, coordinate expansion gives
\[
 \sum_j\chi_{g,j}(z)g(e_j)=|\mathcal M|^{-1}z;
\]
summing over $g$ proves exact reconstruction. The weighted bounds used
below are stronger than the unweighted reconstruction bound. For
$h(z)=\|\operatorname{core}z\|_1$ or
$h(z)=\|\operatorname{vowel}_L(\operatorname{core}z)\|_1$, coordinate
expansion gives
\[
 \sum_{g,j}|\chi_{g,j}(z)|h(e_j)
 =|\mathcal M|^{-1}\sum_g h(g^{-1}z).
\]
For each fixed motion, the weighted coordinate sum is exactly
$h(g^{-1}z)$, by expanding in the frame basis; the core and grade weights
commute with motion. The finite patch-cost estimate gives
$h(g^{-1}z)\le R_Fh(z)$, and therefore yields
\[
 \sum_{g,j}|\chi_{g,j}(z)|h(e_j)\le R_Fh(z).
 \tag{2}\label{eq:ex-motion-weighted}
\]
This weighted identity also explains why one must include the motion cost:
a motion need not preserve the frame's $\ell^1$ norm.
The same coefficient formula with $h(z)=\|z\|_1$ gives the unweighted
budget used for finite-wing additions:
\[
 \sum_{g,j}|\chi_{g,j}(z)|\le R_F\|z\|_1,
\]
by the same finite patch-cost estimate.
Assume each grade weight in $\operatorname{vowel}_L$ lies in $[1,L^d]$;
this is the grading convention used below. The selected purchase inequalities are
\[
 C\ge1,\quad \eta\ge0,\quad
 \eta(1+s)R_F\le\operatorname{barunit}(L),\quad
 (C+\eta s)R_F\le K_{J(m+1)}.
 \tag{1}\label{eq:ex-interpolation-buy}
\]
For each basis vector $e_j$, set $v_j=I^{-1}e_j$. The fine played column
$V(v_j)$ is formatted and integrable, and its remainder is bounded by its
charge. The selected one-column purchase record has three clauses: the
principal insertion when $j\ne0$, every regular remainder atom with its
support and supremum bounds, and every forced atom with its support-mass
bound. The forced clause applies only to nonzero addresses; the keeper
condition places every such address in the funded wing. The unit coordinate
uses the exact vacuum coupon. The density has an absolutely convergent
principal-plus-regular-plus-forced expansion: the regular and forced
majorants are dominated by the coefficient envelope and remainder bound.
For each component, the corresponding purchase clause gives a coupon. The
coupon-merger theorem sums their densities and profiles termwise; Tonelli
applies to the nonnegative majorants, and each scalar cost sum is bounded by
the given purchase budget. This gives the exact basis coupon with lead
$T(L)e_j$. Before simplification its three budgets are
\[
 \eta(\|\operatorname{core}e_j\|_1+V.\operatorname{charge}(v_j)),\quad
 C D(e_j)+\eta V.\operatorname{charge}(v_j),\quad
 t+2V.\operatorname{charge}(v_j).
\]
Each paid column receipt controls its lead after every motion: the transported
coefficient $X_{g,j}$ satisfies
$\|X_{g,j}-T(L)(g e_j)\|_1\le\epsilon\|\operatorname{core}e_j\|_1$
for all $g,j$. This is the all-motion clause in the one-column payment;
no norm-isometry of a motion is used. The heard charge estimate gives
$V.\operatorname{charge}(v_j)\le s\|\operatorname{core}e_j\|_1$;
$\|\operatorname{core}e_j\|_1\le1$ and
$D(e_j)\ge\|\operatorname{core}e_j\|_1$. Hence the budgets are at most
$\epsilon\|\operatorname{core}e_j\|_1$, $cD(e_j)$, and $t+2s$.
The marked coefficient/profile payment is proved in
Section~\ref{sec:projection-payment-producer}. The forced component is
paid there after the forced allowance is at most one. The selected primary
and regular-card comparisons remain explicit analytic inputs; the
one-scale local transition estimate of Section~\ref{sec:local-transfer}
does not itself show that the selected history satisfies those two
purchase clauses. Every component here is evaluated under one
background law and one transition kernel.

\begin{lemma}[The paid one-step cone used in the interpolation]
\label{lem:ex-selected-cone-step}
Let $K\ge1$, $s\ge1$, $L\ge1$, and put $q=L^{1/8}$,
$e=q^{-3}$. Suppose
\[
  100000D^6(1+s)^2\le q,\quad K\le D,
  \quad64D^2\le c_-,\quad64D^2\le c_+.
\]
Let $B$ be an invertible linear map with both $B$ and $B^{-1}$ of
$\ell^1$ norm at most $K$, commuting with deletion of the unit coordinate.
For vectors $u,v$ set $y=\operatorname{vowel}_L(v)$ and assume
\[
 \|u-By\|_1\le e(\|\operatorname{core}v\|_1+t),
 \qquad r\le K\|\operatorname{core}y\|_1+et,
 \qquad 0\le t\le s\|\operatorname{core}v\|_1.
\]
Then
\[
 r\le c_-\|\operatorname{core}u\|_1,\qquad
 \|v\|_1\le c_+\|u\|_1,\qquad
 \|u\|_1\le c_+L^d\|v\|_1.
\]
\end{lemma}

\begin{proof}
The bank inequality, $K\le D$, and $K,D\ge1$ imply
\[
2Ke(1+s)=\frac{2K(1+s)}{q^3}\le1.
\]
Indeed the assumed lower bound on $q$ is at least $2K(1+s)$, and
$q\ge1$. The error bound and $t\le s\|\operatorname{core}v\|_1$
therefore give
\[
2K\|u-By\|_1\le\|\operatorname{core}v\|_1
 \le\|\operatorname{core}y\|_1,
\]
where the second inequality is the coordinatewise fact that all
grade-weights in $\operatorname{vowel}_L$ are at least one. The elementary
perturbation estimate for $B$ now gives
\[
\|y\|_1\le2K\|u\|_1,\qquad
\|u\|_1\le2K\|y\|_1.
\]
Apply the same estimate after deleting the unit coordinate. The map $B$
commutes with that deletion and deletion is contractive, so the same error
bound holds for $\operatorname{core}u-B(\operatorname{core}y)$. Thus
\[
\|\operatorname{core}y\|_1\le2K\|\operatorname{core}u\|_1.
\]
Also
\[
 et\le e(1+s)\|\operatorname{core}v\|_1
    \le \|\operatorname{core}y\|_1,
\]
since $2Ke(1+s)\le1$ and $2K\ge1$. Hence
\[
 r\le2K\|\operatorname{core}y\|_1
    \le4K^2\|\operatorname{core}u\|_1
    \le c_-\|\operatorname{core}u\|_1.
\]
Finally, the grade weights also satisfy
$\|y\|_1\le L^d\|v\|_1$. Consequently
\[
 \|v\|_1\le\|y\|_1\le2K\|u\|_1\le c_+\|u\|_1,
 \quad
 \|u\|_1\le2K\|y\|_1\le2KL^d\|v\|_1
                         \le c_+L^d\|v\|_1.
\]
The two room assumptions imply the needed $4K^2\le c_-$ and
$2K\le c_+$.
\end{proof}

\begin{proposition}[Common linear interpolation of the purchased columns]
\label{prop:ex-selected-interpolation}
Assume $s\ge1$, $L\ge1$, and $t\ge0$. Write the heading as
$T(L)=B\circ\operatorname{vowel}_L$, where $B$ commutes with core deletion
and $\|B\|_{1},\|B^{-1}\|_{1}\le D_0:=K_{J(m+1)}$. Assume the bank and
room bounds
\[
100000D_0^6(1+s)^2\le L^{1/8},\qquad
64D_0^2\le c_-,\quad 64D_0^2\le c_+,
\]
with $\operatorname{barunit}(L)=L^{-3/8}$ and the purchase inequalities
\eqref{eq:ex-interpolation-buy}. These are exactly the scalar hypotheses of
Lemma~\ref{lem:ex-selected-cone-step}, with $K=D=D_0$.
The basis
coupons have their stated exact density/profile identities and a common
gauge door; their wings are the motion images of the common funded wing.
There is a coarser
voice $Q$ with:
\begin{enumerate}
\item every $Q(v)$ is formatted and has remainder bound
      $a(v)=cR_FD(Iv)$ at the coarse forced allowance;
\item its readout $I_Q(v)=Q(v).x$ is a linear isomorphism;
\item at the centered anchor, the exact transition identity
      $Q(v).\mathrm{density}=\mathcal T_m(V(v).\mathrm{density})$ holds;
\item its lead discrepancy satisfies
\[
 \|Q(v).x-T(L)I(v)\|_1
       \le\epsilon R_F\|\operatorname{core}I(v)\|_1;
 \]
\item the finite-wing witness for the exponential entry and forced profile
      of $V(v)$ has budgets
\[
 R_F\|I(v)\|_1(t+2s)\quad\text{and}\quad a(v),
\]
      respectively.
\end{enumerate}
The coefficient, regular, and forced fields are complex-linear in $v$.
The density and readout obey motion covariance, while the forced profile and
exponential entry obey the corresponding covariance after reindexing their
address dictionaries; the conjugation flag is independent of orientation.
\end{proposition}

\begin{proof}
For each basis coordinate put $v_j=I^{-1}e_j$. Its played coefficient has
$x$-coordinate $e_j$ by the heard identity. Let $Z_j$ be the coefficient
from its elementary coupon. For each motion $g$, the transition slice has a
move taking the coarse origin to the required fine anchor. Reorient $Z_j$
by that move and, on the conjugating branch, apply complex conjugation to all
coefficient fields. The centered transport identities give a coefficient
$Z_{g,j}$ with
\[
 Z_{g,j}.x=g(Z_j.x),\qquad
 Z_{g,j}.\mathrm{density}=g\!\left(\mathcal T_m
                   (V(v_j).\mathrm{density})\right),
\]
with the same remainder and two Tale budgets. The remainder assertion
follows from the turned tab estimate and the coupon's future bound; the
finite wing is carried to its image under the move. Soundness of the
transition door gives the common transport, and the side condition $2\leq
n_m$ makes every move admissible on the selected schedule. This explains the
coordinate, format/remainder, density, and wing transport used below.

Index these transported coefficients by the finite set
\[
 \mathcal I=\{(g,j):g\text{ a signed motion},\ j\text{ a frame coordinate}\}.
\]
For a vector $v$, assign weight
\[
       w_v(g,j)=\chi_{g,j}(Iv).
\]
The finite orbit coefficients are the averaging/rebasing coefficients; their
identity is the finite reconstruction formula
\[
       Iv=\sum_{(g,j)\in\mathcal I}
                w_v(g,j)\,g(e_j).
       \tag{2}\label{eq:ex-interpolation-chips}
\]
Define $Q(v)$ by summing the transported coefficients with these weights,
using the disjoint-address merge that retains their support and label data.
Equation~\eqref{eq:ex-interpolation-chips} makes the operation complex-linear
in $v$. Since the index set is finite, there is no limiting argument in this
assembly.

The $x$-coordinate is the weighted sum of the transported column
coordinates. Its leading part is reconstructed by the finite motion average,
since $T(L)$ intertwines the motion action. For the error use the paid
all-motion estimate and \eqref{eq:ex-motion-weighted} with
$h(z)=\|\operatorname{core}z\|_1$:
\[
\begin{aligned}
 \|Q(v).x-T(L)Iv\|_1
 &\le \sum_{g,j}|\chi_{g,j}(Iv)|\,
       \|X_{g,j}-T(L)(g e_j)\|_1\\
 &\le \epsilon\sum_{g,j}|\chi_{g,j}(Iv)|
                      \|\operatorname{core}e_j\|_1\\
 &\le \epsilon R_F\|\operatorname{core}Iv\|_1.
\end{aligned}
\]

Each transported remainder has price $cD(e_j)$. Summing its
absolute coefficient-weighted price and using
\eqref{eq:ex-motion-weighted} with
$h(z)=\|\operatorname{vowel}_L(\operatorname{core}z)\|_1$ gives
\[
 \sum_{g,j}|\chi_{g,j}(Iv)|cD(e_j)
       \le cR_FD(Iv).
\]
The finite coefficient merge preserves formatting and adds these remainder
budgets, proving item 1.

For the density, linearity of the transition integral is justified by the
integrability of each played column. Motion covariance of the heard density
identifies the transported coupon density with the corresponding motion of
the transitioned basis density. Sum with the weights in (2): the result is
$\mathcal T$ applied to the density of $V(v)$. The regular and forced address arrays are finite coefficient sums, hence
complex-linear in $v$; their address dictionaries are reindexed by motion
and need not themselves be invariant. The heard density is motion-covariant
by hypothesis. The forced profile covariance is an explicit assumption. For the
exponential entry, use the coefficient identity
$\mathrm{density}(U)=\mathrm{eval}(U)\,\mathrm{exponentialEntry}(U)
+\mathrm{forcedDress}(U)$. The profile covariance and lawful dressing
identity give covariance of the last term; subtract it from the density
covariance identity and divide by the nonzero local evaluation. This proves
the needed exponential-entry covariance.
Thus the readout and the two profile outputs obey the finite motion
reindexing identities, without asserting invariance of raw address arrays.
Every summand is supported on the same common stencil, so the output support
is independent of $v$. The unit identity follows because the core of the unit coordinate
vanishes, so the lead error is zero, and the heading preserves the unit;
the scale map preserves the unit normalization.

For each $v$, put $u=Q(v).x$ and $y=\operatorname{vowel}_L(Iv)$.
The heading map gives the two-sided map bound with $K=K_{J(m+1)}$;
(1) bounds the leading error by
$\operatorname{barunit}(L)\|\operatorname{core}(Iv)\|_1$ and the
new charge $a(v)=cR_FD(Iv)$ by $K\|\operatorname{core}y\|_1$.
Heard gives $V.\operatorname{charge}(v)\le s\|\operatorname{core}(Iv)\|_1$.
Apply Lemma~\ref{lem:ex-selected-cone-step} with the room and bank bounds.
It yields the new charge cone and the forward/backward coordinate
estimates. In particular, if $Q(v).x=0$, its estimate
$\|Iv\|_1\le c_+\|Q(v).x\|_1$ gives $Iv=0$, and therefore $v=0$.
Thus the linear coordinate map of the new voice is injective and, in finite
dimension, is an isomorphism.

Finally, transport each finite-wing witness to its motion address, multiply
it by $w_v(g,j)$, and take their finite union. The finite repainting identity and motion covariance show that the exponential entries and forced profiles sum
to those of $V(v)$. The first budget is bounded by
\[
 \sum_{(g,j)}|w_v(g,j)|(t+2s)
       \le R_F\|Iv\|_1(t+2s),
\]
while the second is the already proved remainder-price sum $a(v)$. The
finite-wing addition rule therefore gives the required Tale witness. This
proves item 5 and completes the common-pencil construction.
\end{proof}

The proposition is the finite assembly used by the selected one-scale
transition. The source status of its three purchase clauses is not uniform.
The forced clause is supplied by the large-projection estimate after the
forced allowance is at most one and the keeper contains every nonzero slot.
The regular clause still requires the regular-card estimate for the
small-support, vanishing-jet principal term. The main clause for each
nonunit basis coordinate still requires the selected primary-card and
budget-card comparison. Neither missing clause follows from the finite
pencil. Once the exact component receipts are supplied, their absolutely
convergent density decomposition is merged by the dominated-series argument
above, and the finite-wing merger sums the two budgets using the same motion
coefficients. No compactness or diagonal extraction is involved.

\section{Same-run Cauchy input for the full mesh}
\label{ex-public-same-run-cauchy}

We record the actual Cauchy route used to upgrade the positioned limit to the
whole mesh. This is a same-run statement. Fix a selected run and its campaign,
Winter witness $H$, and Pillars witness $W$, together with the chosen growth
function $f_i$ and the selected history tables $C_i$. For a fixed period $p$,
word length $n$, list of local labels $v$, and physical tuple $x$ which is
both non-hit and an ancestor point, write
\[
 o_j(x)=\operatorname{wrap}_{a_j}\bigl(\operatorname{address}_{a_j}(x)\bigr),
 \qquad
 U_i(x)=C_i.\operatorname{numeric}\bigl(f_i,p,v,o_{f_i}(x)\bigr).
\]
The selected source route proves that $(U_i(x))_i$ is Cauchy from the
growth, price, transfer, and geometric receipts of one fixed run. This
subsection translates the Cauchy implication and the analytic interpolation
step; the upstream proof that the selected run supplies all the receipts is
recorded separately below.

The Cauchy step uses one quantitative insertion comparison. We state its
exact logical form so the upstream estimate is visible. Fix a finite label
set $A$, a reference rung $k$, a later rung $j$, and a terminal mesh index
$i$. Let $U_i$ be the full-volume pilot and $V_{j,i}$ its fixed-rung
interpolant. Suppose $F_0<\infty$ is independent of $j$, $G_j\to\infty$,
and, for every sufficiently large fixed $j$, the same-history receipts give
\[
 |U_i-V_{j,i}|\leq \frac{F_0}{G_j/2}
 \quad\text{for all sufficiently large }i,
 \tag{3}\label{eq:ex-fixed-rung-comparison}
\]
and $(V_{j,i})_i$ is Cauchy at each such fixed $j$. A sufficient exact
condition for this fixed-rung Cauchyness is that its finite mode integrand is
a polynomial in finitely many measurable coefficient fields that converge
pointwise, these mode integrands have a common integrable dominator, and
their positive partition normalizers converge to a positive limit. Dominated
convergence then gives convergence of the numerator integrals, and division
by the normalizers gives convergence of $V_{j,i}$. Under these hypotheses
$(U_i)_i$ is Cauchy: given $\varepsilon>0$, first choose $j$ so that
$F_0/(G_j/2)<\varepsilon/4$, then take the maximum of the threshold in
\eqref{eq:ex-fixed-rung-comparison} and the fixed-$j$ Cauchy threshold for
error $\varepsilon/4$. For later $i,i'$,
\[
 |U_i-U_{i'}|
 \leq |U_i-V_{j,i}|+|V_{j,i}-V_{j,i'}|+|V_{j,i'}-U_{i'}|
 <3\varepsilon/4<\varepsilon.
\]

\begin{lemma}[Holomorphic insertion estimate]
\label{lem:ex-holomorphic-insertion}
Let $S>2$ and let $\Phi$ be holomorphic on $B(0,S)$. If
$|\Phi(z)|\leq F_0$ for all $|z|<S$, then
\[
 |\Phi(1)-\Phi(0)|\leq F_0/(S/2).
 \tag{4}\label{eq:ex-holomorphic-insertion}
\]
\end{lemma}
\begin{proof}
Put $r=S/2$. For every $w$ with $|w|\leq1$, the closed disk
$\overline B(w,r)$ lies in $B(0,S)$ since $1+S/2<S$. The Cauchy integral
formula on its boundary gives $|\Phi'(w)|\leq F_0/r$. Integrating
the derivative on the real segment from $0$ to $1$ proves (4).
\end{proof}

\begin{lemma}[Uniform confluence supplies the fixed-rung vault]
\label{lem:ex-uniform-confluence-vault}
Let $I$ be a finite set, let $\mathcal B$ be a finite set, and let $\mathcal X$
be a nonempty set. Suppose $F_{n,i}:\mathcal X\to\mathbb C$ and
$G_{n,i,b}:\mathcal X\to\mathbb C$ have uniform eventual bounds: for each
$i\in I$ there is $c_i\geq0$ such that eventually
$\sup_{x\in\mathcal X}|F_{n,i}(x)|\leq c_i$, and for each $i,b$ there is
$d_{i,b}\geq0$ such that eventually
$\sup_{x\in\mathcal X}|G_{n,i,b}(x)|\leq d_{i,b}$. For any weights
$w_i>0$, there are finite $x_0,y_0$ such that, eventually in $n$,
\[
 w_i|F_{n,i}(x)|\leq x_0,\qquad
 \sum_{b\in\mathcal B}w_i|G_{n,i,b}(x)|\leq y_0
 \quad(i\in I,\ x\in\mathcal X),
\]
where one may take
\[
 x_0=\sum_{i\in I}w_i c_i,
 \qquad y_0=\sum_{i\in I}\sum_{b\in\mathcal B}w_i d_{i,b}.
 \]
\end{lemma}
\begin{proof}
There are only finitely many bounds. Intersect their eventual tails to get one
index $n_0$ after which all hold at once. For each fixed $i$, positivity of
the weights and nonnegativity of every summand give
$w_i|F_{n,i}(x)|\leq w_i c_i\leq x_0$. Likewise
$\sum_b w_i|G_{n,i,b}(x)|\leq\sum_b w_i d_{i,b}\leq y_0$.
This is precisely the two-part vault bound used for the scalar and box-marked
fields. In the selected construction, the source confluence statement gives
an eventual uniform bound over configurations for each fixed rung and each
field. Since the label set and the box set of a fixed finite torus are finite,
their thresholds can be intersected once. No compactness or extraction is
needed.
\end{proof}

The fixed-rung vaults in this argument are produced by the same history, not
by an additional choice of subsequence. At a fixed rung $j$, the retained
history interpolation gives uniform confluence of each scalar insertion and
each box-marked insertion. Applying Lemma~\ref{lem:ex-uniform-confluence-vault}
to the finite label set, with weights $w_i$, gives the constants $t,v$ in the
later-rung vault. Applying it at the opened rung $k$ gives $x,y$ in the
opened-rung vault. The constants may depend on $k,j$, the finite labels, and
the weights, but the opened-rung constants used in the Cauchy comparison are
held fixed as $j$ varies. The positive weights are obtained from the paid
allowances $B_i$ by $w_i=\zeta_i/(|B_i|+1)$; indeed
$w_i|B_i|\leq\zeta_i$ and $w_i>0$ when $\zeta_i>0$.

In the public same-run application, these are actual receipts of the selected
history. Its branch-stability estimate gives uniform confluence at every
fixed rung for the scalar and box-marked fields, and its credit estimate gives
the finite constants $B_i$. The preceding finite sum therefore supplies the
two vaults required at the opened rung and at each later rung. The history's
scale profile supplies the positive weights and their march relation. Its
selected position map supplies the remaining geometry: ancestor-point
separation gives injectivity at the chosen rung, hence occupancy one, and
continues to give the required Drake inequalities at later rungs. These
properties are intersected with the eventual pond, positive-reserve, and
delay tails before fixing the rung; thus the estimates are simultaneous for
the finite label family and hold along the original volume sequence.

To apply Lemma~\ref{lem:ex-holomorphic-insertion}, let $P_z(U)$ be the
selected census-mode coefficient of the restored field at rung $j$, with its
principal and forced field families multiplied by $z$, and put
\[
 \Phi(z)=Z_{N,p}^{-1}\int P_z(U)\,d\mu_j(U).
\]
The finite restored polynomial depends polynomially on $z$. The quantitative
majorant establishing integrability of its coefficients is given below.

Here is the selected-field verification of its uniform bound, conditional on
the served-slab and uniform price receipts at rung $j$. Write $B_j$ for its
Drake separation scale and $E_j$ for the common support exponent. Assume
$B_j\geq8$, $E_j\geq2$, and $S_j=e^{B_j/64}>2$. The last inequality is a
separate late-scale condition; it does not follow from $B_j\geq8$ alone.
For a nonzero polynomial component on a support $X$, each nonzero mode has at
least two markers, and the support-placement condition puts every named
origin of those markers in $X$. Drake separation says distinct origins of
that mode are at max-distance at least $B_j^2$.
The connected-halo tour estimate in Section~\ref{ex-centered-construction}
therefore gives
\[
 B_j\,|I|\leq64\operatorname{cover}(B_j,X),
 \qquad \operatorname{cover}(B_j,X)\geq B_j/64,
\]
for its marker set $I$. Thus $|z|<e^{B_j/64}\leq
e^{\operatorname{cover}(B_j,X)}$. Scaling the component by $z$ multiplies
its envelope by $|z|$. Its original weighted support bill contains the
factor $e^{E_j\operatorname{cover}(B_j,X)}$; since
$|z|\leq e^{2\operatorname{cover}(B_j,X)}$, the same bill pays the scaled
envelope at exponent $E_j-2$. This argument applies to every active support
in both the principal and forced families; zero components cost nothing.
The correction families are scaled, while the exponential-entry and marked
activity carriers are fixed. The support, gauge, and minimum-degree
conditions remain true under scalar multiplication; the envelopes in the
support ledger scale by $|z|$. Hence the same served decomposition with those
scaled correction families witnesses servedness of the restored field at
exponent $E_j-2$. This does not assert that the restore operation itself is
linear. For integrability, the vault bounds the finitely many scalar and
marked fields and the served ledger bounds the finitely many correction
coefficients. The selected lawfulness upper bound for the action and the
vacuum-activity envelope give the pointwise exponential majorant proved in
Section~\ref{ex-centered-steer}; all coefficient fields are measurable.
Expanding the finite restored polynomial therefore gives a measurable,
integrable majorant for each coefficient of $P_z$, uniform for $z$ in the
disk. The prior is a probability measure. Thus the integral defining $\Phi$
is holomorphic on $B(0,S_j)$.

To pass to the opened rung, the served witness at the lowered exponent is
propagated by the downward servedness induction: at each intervening scale
the retained matchup and wake data give the transition hypotheses, and the
exact push identity fixes each marked coefficient. The induction's
integrability step follows from the transition's finite probability mass.
This is the served-transfer result of
Section~\ref{ex-centered-construction}, with its displayed one-step response
and payment hypotheses; Fubini supplies the mode-integral identity below,
but by itself does not imply servedness. The backward transition from rung
$j$ to $k$ is finite, and at each transition, if $K_m(U,V)$ denotes its
kernel from finer configuration $U$ to coarser configuration $V$, its
normalization is $\int K_m(U,V)\,d\mu_m(V)=1$ for every $U$. For each
integrable coefficient $q$, the push identity is
\[
 q_m(V)=\int K_m(U,V)q_{m+1}(U)\,d\mu_{m+1}(U).
\]
Moreover, nonnegativity and normalization give
\begin{align*}
 &\int\!\int K_m(U,V)|q_{m+1}(U)|\,d\mu_{m+1}(U)d\mu_m(V)\\
 &\hspace{2em}=\int |q_{m+1}(U)|\,d\mu_{m+1}(U)<\infty.
\end{align*}
Fubini therefore proves
\[
 \int q_m(V)\,d\mu_m(V)
 =\int q_{m+1}(U)\left(\int K_m(U,V)\,d\mu_m(V)\right)d\mu_{m+1}(U)
 =\int q_{m+1}(U)\,d\mu_{m+1}(U).
\]
This proves the mode-integral preservation at each transition. The opened-rung vaults and
finite-volume mode estimate of Section~\ref{ex-centered-sail} then give
exactly
\[
 |\int P_z\,d\mu_j|\,\operatorname{toll}_w(A)
 \leq Z_{N,p}\exp(|\mathrm{Site}_k|\mathcal T_k+|A|(x+y+2)),
\]
\[
 |\Phi(z)|\leq F_0:=
 \frac{\exp(|\mathrm{Site}_k|\mathcal T_k+|A|(x+y+2))}
      {\operatorname{toll}_w(A)}.
\]
All weights, the opened rung, and its vault budgets are fixed as $j$ varies,
so $F_0$ is independent of $j$. This checks the bound hypothesis of the
holomorphic estimate from the explicit scaling, backward transfer, and
integrated-mode receipts.

In the public pilot comparison, the finite transition recurrence identifies
$U_i=\Phi(1)$ and $V_{j,i}=\Phi(0)$, with the same terminal law and
normalization. The terminal birth identity supplies the value at $z=1$;
the later-rung restoration with its correction families set to zero supplies
the value at $z=0$. The opened rung $k$ is used to price this value after
transfer. The transition
preserves each mode integral, so this identification introduces no change of
partition function. Thus (4) is precisely
\eqref{eq:ex-fixed-rung-comparison}, not an additional compactness input.

There is an exact finite identity connecting this pilot observable to the
ordinary finite-volume Wilson observable. For the census mode containing all
$n$ labels, all copy switches are the identity. The product of the local
insertions is therefore exactly the finite-volume numeric observable, at the
same address tuple $o_{f_i}(x)$. Thus Cauchyness of the pilot sequence is
Cauchyness of $U_i(x)$ itself; this identity introduces no error term and no
subsequence.

To pass from ancestor points to the full positional function, use that
ancestor points are dense. The full-mesh local Lipschitz receipt on a compact
subset of the non-hit set, together with $a_i\to0$, gives asymptotic
uniform equicontinuity there: if the continuum points are within $\delta/2$,
their lattice addresses are within $\delta$ once $2c_n a_i<\delta/2$.
For compact $K\subset U=\mathrm{NoHit}$, choose a compact neighborhood
$K\Subset K\prime\Subset U$ and a common asymptotic modulus on $K\prime$.
A finite $\delta/3$-net for $K$ can be chosen with centers in the dense
ancestor set intersected with $U$. Pointwise Cauchyness at these finitely
many centers gives one threshold. For any two later indices and any $x\in K$,
compare their values at $x$ to the value at a net center: the two modulus
errors are each below $\varepsilon/3$, and the center values differ by
less than $\varepsilon/3$. Hence the functions are uniformly Cauchy on $K$;
completeness gives a uniform limit, and the same modulus makes it continuous.
The generic finite-difference bounds in Section~\ref{sec:ex-fullmesh-fixedbounds}
provide the required bounds and local Lipschitz receipts from the same run.
The application is conditional on the following same-run receipts: uniform
branch confluence and a paid credit for each fixed finite label set; the
pond, migration, occupancy, and Drake data needed to construct the transition
fields; a served witness at each chosen later rung with its displayed
support-price bounds; the finite downward service-transfer hypotheses; and a
fixed-rung Cauchy sequence. It also uses the full-mesh local Lipschitz bound
on each compact subset. Given these receipts, (3) is the quantitative
insertion estimate just proved, and the pointwise Cauchy and compact-uniform
limits required by the full-mesh positioned construction follow for every
period, length, and label word.
Together with the same-run finite-difference bounds in
Section~\ref{sec:ex-fullmesh-fixedbounds}, these receipts yield the smooth
positioned correlation along the full mesh, without further extraction.
For source correspondence, the branch-stability result at a fixed rung gives
confluence of the scalar insertion and of every box-marked insertion; the
credit result gives the paid amounts. The finite-sum lemma above converts
those uniform bounds into the two vault inequalities. At the opened rung the
positive weights are fixed, and at a later rung the marched weights are used
with later-rung confluence. The pond package is passed from the selected
run's scale-field package, and its history recurrence builds the per-label
transition fields used by the served-slab step. These source fields are
mapped in the evidence record; the mathematical production of them is not
claimed here. The detailed local service propagation and its explicit
nonlinear response input are stated in Section~\ref{ex-centered-construction}.
The alternative detector route through trained rows and null-row sweeps is
not called by this public construction.
